\chapter*{Úvod}
\phantomsection
\addcontentsline{toc}{chapter}{Úvod}

Semestrální práce se zabývá návrhem a implementací palubního počítače pro vozidla bez centrální datové sběrnice. Konkrétním vozidlem, ve kterém bude počítač instalován, je Škoda Felicia 1.3 MPI.

V současné době je palubní počítač standardem v nových vozidlech. Díky rozmachu elektromobility a vývoji počítačových technologií se palubní počítače s informacemi o provozu vozidla staly běžnou výbavou.
Starší vozidla, jako třeba Škoda Felicia, tuto možnost měla pouze jako příplatkovou. Vzhledem k tehdejší cenovce palubního počítače APRI TC-6 (případně TC-P6 nebo TM-3 pro dieselové motory) jím bylo vybaveno pouze minimum vozidel a dnes jsou spíše raritou.

Jelikož se starší vozidla těší stále velké popularitě v komunitě nadšenců, rozhodl jsem se takový palubní počítač navrhnout a implementovat.

Základním kamenem pro palubní počítač ve starším voze je mikrokontrolér, který bude zajišťovat sběr dat a také jejich zpracování a prezentaci uživateli. 
V nových vozech je pro toto přizpůsoben řídící systém motoru, který veškerá data ze snímačů sbírá a následně je poskytuje pro zpracovaní přes CAN.\footnote{CAN - centrální datová sběrnice (controller area network)}

Cílem práce je navrhnout a realizovat hardwarové měřicí obvody pro sběr telemetrických dat z vozidla, konkrétně pro snímání rychlosti, otáček motoru, teploty okolí vozidla, provozních teplot motoru a dalších parametrů.

Následující text je rozdělen na dvě hlavní části – teoretickou a praktickou. Teoretická část se věnuje nezbytné rešerši hardwarových komponent, rozboru problematiky palubních systémů obecně a bezpečnosti uživatele při užívání počítače v provozu.
Praktická část následně popisuje samotnou realizaci projektu; zahrnuje návrh měřicích obvodů, tvorbu uživatelského rozhraní na základě externího grafického návrhu,
programovou implementaci a závěrečné testování celého řešení v reálném provozu vozidla.



% Úvod studentské práce, např\,\dots

% Nečíslovaná kapitola Úvod obsahuje \uv{seznámení} čtenáře s~problematikou práce.
% Typicky se zde uvádí:
% (a) do jaké tematické oblasti práce spadá, (b) co jsou hlavní cíle celé práce a (c) jakým způsobem jich bylo dosaženo.
% Úvod zpravidla nepřesahuje jednu stranu.
% Poslední odstavec Úvodu standardně představuje základní strukturu celého dokumentu.

% Tato práce se věnuje oblasti \acs{DSP} (\acl{DSP}), zejména jevům, které nastanou při nedodržení Nyquistovy podmínky pro \ac{symfvz}.%
% \footnote{Tato věta je pouze ukázkou použití příkazů pro sazbu zkratek.}

% Šablona je nastavena na \emph{dvoustranný tisk}.
% Nebuďte překvapeni, že ve vzniklém PDF jsou volné stránky.
% %Je to proto, aby důležité stránky jako např.\ začátky kapitol začínaly po vytisknutí a svázání vždy na pravé straně.
% %
% Pokud máte nějaký závažný důvod sázet (a~zejména tisknout) jednostranně, nezapomeňte si přepnout volbu \texttt{twoside} na \texttt{oneside}!